
\begin{obeactivity}
\textbf{Aktivitas 15.1: Eksperimen Efek Longsoran (\textit{Avalanche Effect}) pada Hash}
1. Gunakan alat bantu online (seperti CyberChef) untuk menghitung SHA-256 dari kalimat: \texttt{"Kriptografi itu Menyenangkan"}.
2. Ubah satu karakter kecil, misalnya menjadi: \texttt{"Kriptografi itu menyenangkan"} (huruf 'm' kecil).
3. Bandingkan kedua hasil hash tersebut.
4. Amati berapa banyak bit yang berubah dan diskusikan mengapa properti ini sangat penting bagi keamanan fungsi hash.
\end{obeactivity}

Langkah-langkah pada Aktivitas 15.1 dapat dijalankan langsung melalui program
pada Listing~\ref{lst:hash-avalanche}. Program itu menghitung SHA-256 dari
kedua kalimat, membandingkan digestnya, dan menghitung berapa bit yang berbeda.
Hasil yang seharusnya Anda peroleh diperlihatkan pada
Gambar~\ref{fig:longsoran-sha256}: $115$ dari $256$ bit berubah, yaitu
$44{,}92\%$ --- mendekati setengah, sebagaimana disyaratkan oleh uji longsoran
pada Subbagian~\ref{subsec:anatomi-sha256}. Bila angka yang Anda peroleh
berbeda jauh dari setengah, ada yang salah pada pesan yang di-hash, bukan pada
SHA-256.

\lstinputlisting[
    language=Python,
    caption={Efek longsoran SHA-256: perbandingan bit dua digest},
    label={lst:hash-avalanche}
]{\codefile{python/hash_avalanche.py}}

Program yang sama tersedia dalam C++ (\texttt{code/cpp/hash\_avalanche.cpp})
dan MATLAB (\texttt{code/matlab/hash\_avalanche.m}); keduanya memuat
implementasi SHA-256 sendiri karena bahasanya tidak menyediakannya secara
bawaan.

\begin{obeactivity}
\textbf{Aktivitas 15.2: Menelusuri Jejak Kompresi SHA-256}
\label{act:hash-sha256}
1. Jalankan \texttt{python code/python/hash\_uji.py} dan catat keadaan register
   $a$ dan $e$ pada putaran 1, 2, 16, 32, 48, dan 64.
2. Bandingkan digest yang dihasilkan dengan keluaran \texttt{hashlib.sha256}.
3. Ubah pesan masukan menjadi \texttt{abd} dan amati pada putaran ke berapa
   keadaan register mulai berbeda dari \texttt{abc}.
4. Periksa hasil perhitungan padding untuk tiga panjang pesan yang tercetak:
   mengapa panjang 448 bit justru memerlukan dua blok, bukan satu?
\end{obeactivity}

Berbeda dengan Aktivitas 15.1 yang memakai \texttt{hashlib} sebagai kotak hitam,
aktivitas ini membuka isi kotaknya. SHA-256 ditulis dari definisi sehingga setiap
putaran dapat diperiksa satu per satu. Bagian jadwal pesan dan fungsi kompresi
(Persamaan~\eqref{eq:sha256-jadwal} dan~\eqref{eq:sha256-kompresi}) diperlihatkan
pada Listing~\ref{lst:hash-jadwal}, sedangkan aturan padding
(Persamaan~\eqref{eq:padding-sha}) beserta pengulangan bloknya ada pada
Listing~\ref{lst:hash-padding}.

\lstinputlisting[
    language=Python,
    caption={Jadwal pesan dan fungsi kompresi SHA-256 selama 64 putaran},
    label={lst:hash-jadwal},
    firstline=60,
    lastline=90
]{\codefile{python/hash_uji.py}}

\lstinputlisting[
    language=Python,
    caption={Aturan padding Merkle--Damgård dan pengulangan blok SHA-256},
    label={lst:hash-padding},
    firstline=92,
    lastline=110
]{\codefile{python/hash_uji.py}}

Hasil menjalankan program untuk pesan \texttt{abc} memperlihatkan bahwa sesudah
satu putaran saja register $a$ sudah berubah dari \code{6a09e667} menjadi
\code{5d6aebcd}, dan sesudah 64 putaran digestnya berakhir dengan
\code{ba7816bf} --- tepat sama dengan nilai baku. Perhatikan juga bahwa peredaran
nilai pada register tidak memperlihatkan pola apa pun terhadap kesederhanaan
masukan \texttt{abc}. Itulah wujud konkret sifat \emph{diffusion} yang dibahas
pada Subbagian~\ref{subsec:konstruksi-hash}.

\begin{obeactivity}
\textbf{Aktivitas 15.3: Membuktikan Serangan Panjang-Ekstensi}
1. Jalankan program dan catat nilai $H(K \| m)$ untuk $K = \texttt{"kunci"}$ dan
   $m = \texttt{"pesan"}$.
2. Perhatikan bahwa penyerang hanya mengetahui panjang $|K \| m| = 10$ byte,
   bukan isi $K$.
3. Bandingkan tag palsu yang dihitung dari keadaan internal dengan hash yang
   dihitung \texttt{hashlib} atas pesan yang diperpanjang. Apakah keduanya sama?
4. Diskusikan: berapa operasi tambahan yang dibutuhkan penyerang dibandingkan
   dengan menghitung hash pesan biasa?
\end{obeactivity}

Serangan pada Aktivitas 15.3 memperlihatkan bahwa skema $H(K \| m)$ dapat
ditembus tanpa mengetahui kunci sama sekali. Seluruh serangan itu dijalankan oleh
fungsi pada Listing~\ref{lst:hash-ekstensi}, yang hanya melanjutkan kompresi dari
digest yang sudah diketahui. Hasil numeriknya dibahas pada
Contoh~\ref{ex:length-extension}.

\lstinputlisting[
    language=Python,
    caption={Inti serangan panjang-ekstensi: melanjutkan kompresi dari digest},
    label={lst:hash-ekstensi},
    firstline=117,
    lastline=130
]{\codefile{python/hash_uji.py}}

\begin{obeactivity}
\textbf{Aktivitas 15.4: Menguji MAC: HMAC, Vektor RFC 4231, dan Waktu-Konstan}
1. Jalankan program dan periksa apakah kelima vektor uji HMAC-SHA-256 cocok
   dengan keluaran \texttt{hmac.new}.
2. Perhatikan TC6: kuncinya 131 byte, melebihi ukuran blok 64 byte. Cabang mana
   pada Persamaan~\eqref{eq:hmac-panjang-kunci} yang dipakai?
3. Hitung sendiri batas ulang tahun untuk $n = 128$, $160$, $256$, dan $512$,
   lalu bandingkan dengan nilai yang tercetak.
4. Pada bagian terakhir, kedua fungsi pemeriksa tag menolak tag yang diubah.
   Jelaskan mengapa cara pertama tetap berbahaya meskipun jawabannya benar.
\end{obeactivity}

Listing~\ref{lst:hash-hmac} memuat HMAC yang ditulis mengikuti
Persamaan~\eqref{eq:hmac-rumus}, sedangkan Listing~\ref{lst:hash-uji} memuat
program pengujiannya: jejak kompresi, aturan padding, serangan panjang-ekstensi,
kelima vektor RFC 4231, batas ulang tahun, dan perbandingan tag secara
waktu-konstan.

\lstinputlisting[
    language=Python,
    caption={HMAC-SHA-256 dari definisi: penyesuaian panjang kunci dan dua lapis hash},
    label={lst:hash-hmac},
    firstline=136,
    lastline=146
]{\codefile{python/hash_uji.py}}

\lstinputlisting[
    language=Python,
    caption={Program pengujian: jejak kompresi, panjang-ekstensi, vektor RFC 4231,
    batas ulang tahun, dan perbandingan waktu-konstan},
    label={lst:hash-uji},
    firstline=166,
    lastline=263
]{\codefile{python/hash_uji.py}}

Sesudah aktivitas ini, mahasiswa diharapkan dapat menjawab pertanyaan yang
menjadi inti bab ini: bukan sekadar \emph{apa} fungsi hash itu, melainkan
\emph{mengapa} setiap pilihan konstruksinya diambil, dan apa yang terjadi bila
pilihan itu keliru.
